\TopicDivider{02}{Analytical Conjugacy vs Approximate Inference}{Closed-form elegance versus scalable modern sampling algorithms.}

\begin{frame}{Analytical Conjugacy: Beta-Binomial Derivation}
  \MathDerivationStep{1}{Conjugate Kernel}{\theta \sim \text{Beta}(\alpha, \beta) \implies p(\theta) \propto \theta^{\alpha - 1} (1 - \theta)^{\beta - 1}}{%
    The Beta family acts as a conjugate prior for binomial observation processes, closed under Bayesian updating.
  }
  \vspace{.12cm}
  \MathDerivationStep{2}{Exact Posterior Update}{y \mid \theta \sim \text{Binomial}(n, \theta) \implies \theta \mid y \sim \text{Beta}(\alpha + y, \, \beta + n - y)}{%
    Direct kernel multiplication: \(\theta^{(\alpha - 1) + y} (1 - \theta)^{(\beta - 1) + (n - y)}\). Hyperparameters update additively without numerical integration.
  }
\end{frame}

\begin{frame}[fragile]{Production Vectorized Probabilistic Kernel}
  \begin{HighlightedCodeBlock}[numbers=left]{Python}{7,8,9,10}
import numpy as np

def conjugate_beta_binomial_update(alpha_prior: float, beta_prior: float, 
                                   successes: int, trials: int) -> dict:
    """Calculates exact analytical posterior parameters and 95% Credible Interval."""
    failures = trials - successes
    alpha_post = alpha_prior + successes
    beta_post = beta_prior + failures
    
    post_mean = alpha_post / (alpha_post + beta_post)
    post_var = (alpha_post * beta_post) / (((alpha_post + beta_post)**2) * (alpha_post + beta_post + 1))
    return {"alpha": alpha_post, "beta": beta_post, "mean": post_mean, "variance": post_var}
  \end{HighlightedCodeBlock}
  \vspace{.10cm}
  \Citation{Production implementations maintain numerical stability by evaluating logarithms of Beta function ratios \(\ln \text{B}(\alpha, \beta)\).}
\end{frame}

\begin{frame}{Approximation Frontiers: MCMC vs Variational Inference}
  \CompareGrid{Markov Chain Monte Carlo (MCMC)}{%
    \begin{itemize}\setlength{\itemsep}{.25em}
      \item \textbf{Mechanism}: Generates dependent samples from the exact posterior using Markov chains (e.g., NUTS, HMC).
      \item \textbf{Guarantee}: Exact in the asymptotic limit (\(S \to \infty\)).
      \item \textbf{Bottleneck}: High computation time for large datasets or deep models.
    \end{itemize}
    \PillBadge[BayesBlue]{EXACT ASYMPTOTICS}
  }{Variational Inference (VI)}{%
    \begin{itemize}\setlength{\itemsep}{.25em}
      \item \textbf{Mechanism}: Casts inference as an optimization problem minimizing \(\text{KL}(q_\phi(\theta) \parallel p(\theta \mid y))\) via the ELBO.
      \item \textbf{Guarantee}: Approximate density with bounded analytical divergence.
      \item \textbf{Advantage}: Scalable to massive data via stochastic gradients.
    \end{itemize}
    \PillBadge[BayesTeal]{DETERMINISTIC SCALE}
  }
\end{frame}

\begin{frame}{The Computational Complexity Spectrum}
  \StatGridThree{%
    \StatCard{$O(1)$}{Exact Algebra}{Conjugate Analysis}{Closed-form formulas for exponential families}%
  }{%
    \StatCard{$S \to \infty$}{NUTS / HMC}{MCMC Sampling}{Asymptotically exact ergodic convergence}%
  }{%
    \StatCard{$D > 10^7$}{ELBO Gradient}{Variational Inference}{Scalable stochastic gradient descent}%
  }
  \vspace{.15cm}
  \TakeawayBanner{Choose MCMC when exact uncertainty calibration is mandatory; choose VI when data scale forbids iterative sampling.}
\end{frame}
